% Included by manuscript.tex via % Included by manuscript.tex via % Included by manuscript.tex via % Included by manuscript.tex via \input{section_5_3_schemes}.
% Not standalone: no preamble, no \begin{document}.
%
% Numbers are the COMPLETE survey: 120 state points, 4 systems x 6
% polydispersities x 5 volume fractions, every point converged and screened.
% Source data: tools/survey.jsonl. Regenerate the tables and figure with
%     python tools/section53_survey.py --summarise survey.jsonl
%     python tools/section53_survey.py --plot survey.jsonl
%
% CAUTION FOR WHOEVER EDITS THIS. Five successive versions of this section
% have supported five different conclusions, each overturned by adding data.
% Do not tighten the language below without adding state points first.

\subsection{Comparison of the approximation schemes}
\label{sec:schemes}

We computed all six schemes alongside the exact multicomponent result over a
grid of relative polydispersity $s = 0.05$--$0.40$ and volume fraction
$\varphi = 0.05$--$0.40$ for four systems: hard spheres and adhesive hard
spheres ($\tau = 0.3$, $\delta = 0.02\sigma$) in the Percus--Yevick closure,
and a square-\emph{shoulder} fluid ($\varepsilon = +1\,k_BT$,
$\delta = 0.1\sigma$) and a
hard-core Yukawa fluid (screening length $0.5\sigma$, contact strength
$2\,k_BT$) in the hypernetted-chain closure. That is 120 state points, every
one converged and screened for physicality. The error measure is
$\max|I_{\mathrm{approx}}/I_{\mathrm{exact}} - 1|$ over
$0.3 \le Q\sigma \le 20$, with five structure-factor classes and sixty
form-factor classes. The narrow-well systems were computed at the resolution
their well width demands (\S\ref{sec:welltrap}): 300 points per diameter for
the square shoulder, 1500 for the adhesive spheres.

\subsubsection{Three systems share a pattern; the Yukawa fluid does not}

\begin{quote}\small\itshape
Note on the sign convention. The implementation writes the pair potential as
$u(r) = \varepsilon$ for $\sigma < r < \sigma+\delta$ and forms
$\exp(-u)$, so a POSITIVE $\varepsilon$ gives $\exp(-\varepsilon) < 1$ and is
REPULSIVE --- a square shoulder. Attraction requires $\varepsilon < 0$. The
survey was run at $\varepsilon = +1\,k_BT$ and is therefore a shoulder
throughout; the calculations are unaffected, but every description of this
row as a ``square well'' was wrong and has been corrected. Before submission,
consider adding a genuinely attractive row ($\varepsilon < 0$): an attractive
well and a repulsive shoulder are different physics, and the adhesive-sphere
row is the only attractive system currently in the survey.
\end{quote}

For the hard-sphere, square-shoulder and adhesive-sphere fluids the same three
schemes divide the $(s,\varphi)$ plane between them, in a consistent
arrangement (Table~\ref{tab:map}): the scaling approximation at low
polydispersity and low volume fraction, the local monodisperse approximation
through the middle of the plane and increasingly as $s$ grows, and the van der
Waals one-fluid scheme along the high-$\varphi$ edge. The boundaries move
between potentials --- the scaling region shrinks as attraction is added, and
vanishes altogether for the adhesive spheres --- but the arrangement does not.
In all three, the decoupling and partial-structure-factor schemes are never
the best choice at any state point.

The hard-core Yukawa fluid behaves differently in both respects. There the
decoupling approximation is the \emph{best} scheme at 21 of 30 points, and
every scheme carries a median error between 29 and 58\,\%. That is a statement
about which approximation is least bad rather than about any of them being
usable: a parameter extracted through any of these schemes for such a system
carries a systematic error of tens of per cent.

This is not an artefact of the particular Yukawa parameters. Varying the
screening length over a fivefold range and the contact strength by a factor
of two and a half leaves both halves of the statement intact: decoupling is
the best scheme at every state point tested, and the median error across all
six schemes stays between $0.55$ and $0.79$ throughout
(Table~\ref{tab:yukawa}).

\begin{table}[htbp]
\centering
\caption{The Yukawa behaviour is insensitive to the interaction parameters.
$\lambda$ is the screening length in units of $\sigma$ and $K$ the contact
strength in $k_BT$; the median is taken across all six approximation schemes.}
\label{tab:yukawa}
\small
\begin{tabular}{ccccll}
\toprule
$\lambda$ & $K$ & $s$ & $\varphi$ & best scheme & median\\
\midrule
0.5 & 2.0 & 0.20 & 0.20 & decoupling (0.42) & 0.60\\
0.5 & 2.0 & 0.30 & 0.30 & decoupling (0.35) & 0.79\\
0.2 & 2.0 & 0.20 & 0.20 & decoupling (0.18) & 0.55\\
0.2 & 2.0 & 0.30 & 0.30 & decoupling (0.24) & 0.76\\
1.0 & 2.0 & 0.20 & 0.20 & decoupling (0.39) & 0.70\\
0.5 & 5.0 & 0.20 & 0.20 & decoupling (0.44) & 0.73\\
\bottomrule
\end{tabular}
\end{table}

Two further combinations --- the longest-ranged and the strongest attraction,
both at $s = \varphi = 0.30$ --- failed to converge. That is expected on
approaching a spinodal rather than a defect, but it does mean the Yukawa
behaviour is characterised on its convergent side only.

\begin{table}[htbp]
\centering
\caption{Scheme with the smallest maximum relative error at each state point.
Three systems share an arrangement; the Yukawa fluid is the exception.
``dec'' is the decoupling and ``part'' the partial-$S_{ij}$ scheme.}
\label{tab:map}
\small
\begin{tabular}{c|ccccc|ccccc}
\toprule
& \multicolumn{5}{c|}{hard sphere, $\varphi$}
& \multicolumn{5}{c}{square shoulder, $\varphi$}\\
$s$ & 0.05 & 0.10 & 0.20 & 0.30 & 0.40 & 0.05 & 0.10 & 0.20 & 0.30 & 0.40\\
\midrule
0.05 & scal & scal & scal & scal & scal & scal & vdW1 & vdW1 & vdW1 & vdW1\\
0.10 & scal & scal & scal & vdW1 & vdW1 & scal & vdW1 & vdW1 & vdW1 & vdW1\\
0.15 & scal & scal & LMA  & LMA  & vdW1 & scal & LMA  & LMA  & LMA  & vdW1\\
0.20 & scal & LMA  & LMA  & LMA  & vdW1 & LMA  & LMA  & LMA  & LMA  & vdW1\\
0.30 & scal & LMA  & LMA  & LMA  & vdW1 & LMA  & LMA  & LMA  & LMA  & vdW1\\
0.40 & LMA  & LMA  & LMA  & LMA  & vdW1 & LMA  & LMA  & LMA  & LMA  & LMA\\
\midrule
& \multicolumn{5}{c|}{adhesive HS, $\varphi$}
& \multicolumn{5}{c}{Yukawa, $\varphi$}\\
$s$ & 0.05 & 0.10 & 0.20 & 0.30 & 0.40 & 0.05 & 0.10 & 0.20 & 0.30 & 0.40\\
\midrule
0.05 & vdW1 & vdW1 & vdW1 & vdW1 & vdW1 & scal & scal & dec  & dec  & dec\\
0.10 & vdW1 & vdW1 & vdW1 & vdW1 & vdW1 & scal & scal & dec  & dec  & dec\\
0.15 & vdW1 & vdW1 & vdW1 & vdW1 & vdW1 & scal & scal & dec  & dec  & dec\\
0.20 & vdW1 & vdW1 & LMA  & LMA  & vdW1 & dec  & scal & dec  & dec  & dec\\
0.30 & LMA  & LMA  & LMA  & LMA  & vdW1 & dec  & dec  & dec  & dec  & dec\\
0.40 & LMA  & LMA  & LMA  & LMA  & vdW1 & dec  & dec  & dec  & part & part\\
\bottomrule
\end{tabular}
\end{table}

\subsubsection{Typical and worst-case accuracy}

Over all 120 points the local monodisperse approximation and the van der Waals
one-fluid scheme are the best choice about equally often (38 and 37 points
respectively), the scaling approximation at 22 and the decoupling
approximation at 21 --- the latter almost entirely from the Yukawa system. The
partial-structure-factor scheme wins twice. Within-5\,\% rates over all
points are 35.8\,\% for vdW1, 31.7\,\% for both the local monodisperse and
scaling approximations, 16.7\,\% for the monodisperse approximation and
11.7\,\% for decoupling and partial $S_{ij}$.

Typical and worst-case accuracy do not agree, and the choice between them
depends on the application (Table~\ref{tab:schemestats}). For hard spheres the
local monodisperse approximation has the smallest median error but the van der
Waals one-fluid scheme the smallest worst case, by a factor of nearly six.
Reporting only one of these would mislead.

\begin{table}[htbp]
\centering
\caption{Median and worst-case error over the thirty state points of each
system. Best in each column in bold.}
\label{tab:schemestats}
\small
\begin{tabular}{lcccccccc}
\toprule
& \multicolumn{2}{c}{hard sphere} & \multicolumn{2}{c}{square well}
& \multicolumn{2}{c}{adhesive HS} & \multicolumn{2}{c}{Yukawa}\\
Scheme & med. & worst & med. & worst & med. & worst & med. & worst\\
\midrule
monodisperse       & 0.181 & 16.79 & 0.212 & 1.48 & 0.099 & 0.76 & 0.568 & 0.96\\
decoupling         & 0.558 & 286.07 & 0.716 & 28.22 & 0.167 & 5.92 & \textbf{0.287} & \textbf{0.42}\\
local monodisperse & \textbf{0.059} & 17.06 & \textbf{0.071} & 1.50 & 0.068 & 0.76 & 0.566 & 0.95\\
partial $S_{ij}$   & 0.558 & 286.11 & 0.717 & 28.22 & 0.167 & 5.92 & 0.290 & \textbf{0.42}\\
scaling            & 0.072 & 66.42 & 0.129 & 6.36 & 0.059 & 1.71 & 0.417 & 0.79\\
vdW1               & 0.085 & \textbf{11.87} & 0.088 & 2.00 & \textbf{0.044} & \textbf{0.61} & 0.575 & 1.11\\
\bottomrule
\end{tabular}
\end{table}

\subsubsection{The decoupling approximation}

The decoupling and partial-structure-factor schemes track each other closely
throughout, which is expected: both replace the distinct pair correlations by
a single common structure factor. Measured against the best scheme available
at the same state point, the decoupling approximation is typically worse by a
factor of

\begin{center}\small
\begin{tabular}{lcc}
\toprule
System & median ratio & worst\\
\midrule
hard sphere & 13.0 & 48.4\\
square well & 14.8 & 43.2\\
adhesive HS & 4.9 & 10.1\\
Yukawa & 1.0 & 1.6\\
\bottomrule
\end{tabular}
\end{center}

\noindent so its penalty falls monotonically as the interaction becomes softer
and longer ranged, until for the Yukawa fluid it disappears.

This is consistent with the known reason for its failure, and with an
independent result. Gazzillo \emph{et al.} (1999) reached the same conclusion
for polydisperse ionic fluids --- that the extended decoupling approximation
behaves poorly across essentially the whole $Q$ range while the scaling
approximation is accurate over the experimentally accessible range --- and
identified the cause: the decoupling approximation neglects the excluded-volume
conditions $g_{\alpha\beta}(r) = 0$ inside the hard cores, because it
replaces the distinct pair correlation functions by a small number of
effective ones. Where the structure factor is dominated by excluded volume the
approximation that discards excluded-volume detail loses most; where it is
dominated by a long-ranged soft tail, it loses least. Our implementation of
the scaling approximation retains both ingredients Gazzillo identifies as
necessary, the volume prefactor and evaluation of the reference structure
factor at a scaled argument, while our decoupling approximation uses a single
common structure factor, the $\lambda = 1$ limit in their notation.

\begin{quote}\small\itshape
Draft note. This section has now been written five times, on 4, 9, 60, 98 and
120 state points, and each version supported a different conclusion: that the
scheme ordering reverses between potentials (an artefact of the defect
described in \S\ref{sec:validation}); that the local monodisperse
approximation is uniformly best; that three schemes divide the plane; that the
systems fall into three regimes including one where the choice barely matters;
and the account above. The fourth of those was an artefact of having run only
the eight lowest-polydispersity adhesive points, which are indeed insensitive
to the choice --- the completed row is not. The $s = \varphi = 0.40$ corner
has now been checked and IS physical --- the diagonal partial structure
factors stay near $0.97$, the intensity is positive throughout, the packing
fraction reproduces exactly, and $\varphi = 0.40$ with a twofold spread in
diameter is nowhere near jamming, polydispersity raising the close-packing
threshold rather than lowering it. The factor of 286 therefore stands as a
genuine failure of the approximation at an ordinary state point, not a
numerical artefact.
\end{quote}

\subsection{Two cautions for the practitioner}

Both of the following were established while validating this work and are easy
to fall foul of.

\paragraph{Grid resolution is set by the narrowest feature.}
For potentials with a narrow attractive well the discretisation must resolve
the \emph{well}, not the particle: $\gtrsim 30\,\sigma/\delta$ points per
diameter. At the default resolution an adhesive-sphere calculation with
$\delta = 0.02\sigma$ is in error by ninety per cent, and the resulting
structure factor is smooth, plausible and wrong. In an earlier version of this
comparison this made all six schemes appear uniformly poor ($0.72$--$0.80$)
for that system, because each was being compared against a reference that was
itself badly in error.

\paragraph{Convergence flags and residuals are not sufficient.}
The closure equations admit multiple fixed points. At one strongly coupled
state two different accelerated solvers converged to residuals of
$\sim10^{-12}$ and returned $S(Q)=17.0$ and $9.7$ respectively, both with
$\min S(Q) < 0$. Since a structure factor is a variance it cannot be negative,
and solutions are therefore screened for physicality in addition to the
residual test. Separately, the SUNDIALS KINSOL solver was observed to report
success after diverging, because exponent clipping in the closure pins values
at $\sim10^{217}$ and defeats its internal stopping criterion; every solve
consequently recomputes its own residual before the result is accepted.
.
% Not standalone: no preamble, no \begin{document}.
%
% Numbers are the COMPLETE survey: 120 state points, 4 systems x 6
% polydispersities x 5 volume fractions, every point converged and screened.
% Source data: tools/survey.jsonl. Regenerate the tables and figure with
%     python tools/section53_survey.py --summarise survey.jsonl
%     python tools/section53_survey.py --plot survey.jsonl
%
% CAUTION FOR WHOEVER EDITS THIS. Five successive versions of this section
% have supported five different conclusions, each overturned by adding data.
% Do not tighten the language below without adding state points first.

\subsection{Comparison of the approximation schemes}
\label{sec:schemes}

We computed all six schemes alongside the exact multicomponent result over a
grid of relative polydispersity $s = 0.05$--$0.40$ and volume fraction
$\varphi = 0.05$--$0.40$ for four systems: hard spheres and adhesive hard
spheres ($\tau = 0.3$, $\delta = 0.02\sigma$) in the Percus--Yevick closure,
and a square-\emph{shoulder} fluid ($\varepsilon = +1\,k_BT$,
$\delta = 0.1\sigma$) and a
hard-core Yukawa fluid (screening length $0.5\sigma$, contact strength
$2\,k_BT$) in the hypernetted-chain closure. That is 120 state points, every
one converged and screened for physicality. The error measure is
$\max|I_{\mathrm{approx}}/I_{\mathrm{exact}} - 1|$ over
$0.3 \le Q\sigma \le 20$, with five structure-factor classes and sixty
form-factor classes. The narrow-well systems were computed at the resolution
their well width demands (\S\ref{sec:welltrap}): 300 points per diameter for
the square shoulder, 1500 for the adhesive spheres.

\subsubsection{Three systems share a pattern; the Yukawa fluid does not}

\begin{quote}\small\itshape
Note on the sign convention. The implementation writes the pair potential as
$u(r) = \varepsilon$ for $\sigma < r < \sigma+\delta$ and forms
$\exp(-u)$, so a POSITIVE $\varepsilon$ gives $\exp(-\varepsilon) < 1$ and is
REPULSIVE --- a square shoulder. Attraction requires $\varepsilon < 0$. The
survey was run at $\varepsilon = +1\,k_BT$ and is therefore a shoulder
throughout; the calculations are unaffected, but every description of this
row as a ``square well'' was wrong and has been corrected. Before submission,
consider adding a genuinely attractive row ($\varepsilon < 0$): an attractive
well and a repulsive shoulder are different physics, and the adhesive-sphere
row is the only attractive system currently in the survey.
\end{quote}

For the hard-sphere, square-shoulder and adhesive-sphere fluids the same three
schemes divide the $(s,\varphi)$ plane between them, in a consistent
arrangement (Table~\ref{tab:map}): the scaling approximation at low
polydispersity and low volume fraction, the local monodisperse approximation
through the middle of the plane and increasingly as $s$ grows, and the van der
Waals one-fluid scheme along the high-$\varphi$ edge. The boundaries move
between potentials --- the scaling region shrinks as attraction is added, and
vanishes altogether for the adhesive spheres --- but the arrangement does not.
In all three, the decoupling and partial-structure-factor schemes are never
the best choice at any state point.

The hard-core Yukawa fluid behaves differently in both respects. There the
decoupling approximation is the \emph{best} scheme at 21 of 30 points, and
every scheme carries a median error between 29 and 58\,\%. That is a statement
about which approximation is least bad rather than about any of them being
usable: a parameter extracted through any of these schemes for such a system
carries a systematic error of tens of per cent.

This is not an artefact of the particular Yukawa parameters. Varying the
screening length over a fivefold range and the contact strength by a factor
of two and a half leaves both halves of the statement intact: decoupling is
the best scheme at every state point tested, and the median error across all
six schemes stays between $0.55$ and $0.79$ throughout
(Table~\ref{tab:yukawa}).

\begin{table}[htbp]
\centering
\caption{The Yukawa behaviour is insensitive to the interaction parameters.
$\lambda$ is the screening length in units of $\sigma$ and $K$ the contact
strength in $k_BT$; the median is taken across all six approximation schemes.}
\label{tab:yukawa}
\small
\begin{tabular}{ccccll}
\toprule
$\lambda$ & $K$ & $s$ & $\varphi$ & best scheme & median\\
\midrule
0.5 & 2.0 & 0.20 & 0.20 & decoupling (0.42) & 0.60\\
0.5 & 2.0 & 0.30 & 0.30 & decoupling (0.35) & 0.79\\
0.2 & 2.0 & 0.20 & 0.20 & decoupling (0.18) & 0.55\\
0.2 & 2.0 & 0.30 & 0.30 & decoupling (0.24) & 0.76\\
1.0 & 2.0 & 0.20 & 0.20 & decoupling (0.39) & 0.70\\
0.5 & 5.0 & 0.20 & 0.20 & decoupling (0.44) & 0.73\\
\bottomrule
\end{tabular}
\end{table}

Two further combinations --- the longest-ranged and the strongest attraction,
both at $s = \varphi = 0.30$ --- failed to converge. That is expected on
approaching a spinodal rather than a defect, but it does mean the Yukawa
behaviour is characterised on its convergent side only.

\begin{table}[htbp]
\centering
\caption{Scheme with the smallest maximum relative error at each state point.
Three systems share an arrangement; the Yukawa fluid is the exception.
``dec'' is the decoupling and ``part'' the partial-$S_{ij}$ scheme.}
\label{tab:map}
\small
\begin{tabular}{c|ccccc|ccccc}
\toprule
& \multicolumn{5}{c|}{hard sphere, $\varphi$}
& \multicolumn{5}{c}{square shoulder, $\varphi$}\\
$s$ & 0.05 & 0.10 & 0.20 & 0.30 & 0.40 & 0.05 & 0.10 & 0.20 & 0.30 & 0.40\\
\midrule
0.05 & scal & scal & scal & scal & scal & scal & vdW1 & vdW1 & vdW1 & vdW1\\
0.10 & scal & scal & scal & vdW1 & vdW1 & scal & vdW1 & vdW1 & vdW1 & vdW1\\
0.15 & scal & scal & LMA  & LMA  & vdW1 & scal & LMA  & LMA  & LMA  & vdW1\\
0.20 & scal & LMA  & LMA  & LMA  & vdW1 & LMA  & LMA  & LMA  & LMA  & vdW1\\
0.30 & scal & LMA  & LMA  & LMA  & vdW1 & LMA  & LMA  & LMA  & LMA  & vdW1\\
0.40 & LMA  & LMA  & LMA  & LMA  & vdW1 & LMA  & LMA  & LMA  & LMA  & LMA\\
\midrule
& \multicolumn{5}{c|}{adhesive HS, $\varphi$}
& \multicolumn{5}{c}{Yukawa, $\varphi$}\\
$s$ & 0.05 & 0.10 & 0.20 & 0.30 & 0.40 & 0.05 & 0.10 & 0.20 & 0.30 & 0.40\\
\midrule
0.05 & vdW1 & vdW1 & vdW1 & vdW1 & vdW1 & scal & scal & dec  & dec  & dec\\
0.10 & vdW1 & vdW1 & vdW1 & vdW1 & vdW1 & scal & scal & dec  & dec  & dec\\
0.15 & vdW1 & vdW1 & vdW1 & vdW1 & vdW1 & scal & scal & dec  & dec  & dec\\
0.20 & vdW1 & vdW1 & LMA  & LMA  & vdW1 & dec  & scal & dec  & dec  & dec\\
0.30 & LMA  & LMA  & LMA  & LMA  & vdW1 & dec  & dec  & dec  & dec  & dec\\
0.40 & LMA  & LMA  & LMA  & LMA  & vdW1 & dec  & dec  & dec  & part & part\\
\bottomrule
\end{tabular}
\end{table}

\subsubsection{Typical and worst-case accuracy}

Over all 120 points the local monodisperse approximation and the van der Waals
one-fluid scheme are the best choice about equally often (38 and 37 points
respectively), the scaling approximation at 22 and the decoupling
approximation at 21 --- the latter almost entirely from the Yukawa system. The
partial-structure-factor scheme wins twice. Within-5\,\% rates over all
points are 35.8\,\% for vdW1, 31.7\,\% for both the local monodisperse and
scaling approximations, 16.7\,\% for the monodisperse approximation and
11.7\,\% for decoupling and partial $S_{ij}$.

Typical and worst-case accuracy do not agree, and the choice between them
depends on the application (Table~\ref{tab:schemestats}). For hard spheres the
local monodisperse approximation has the smallest median error but the van der
Waals one-fluid scheme the smallest worst case, by a factor of nearly six.
Reporting only one of these would mislead.

\begin{table}[htbp]
\centering
\caption{Median and worst-case error over the thirty state points of each
system. Best in each column in bold.}
\label{tab:schemestats}
\small
\begin{tabular}{lcccccccc}
\toprule
& \multicolumn{2}{c}{hard sphere} & \multicolumn{2}{c}{square well}
& \multicolumn{2}{c}{adhesive HS} & \multicolumn{2}{c}{Yukawa}\\
Scheme & med. & worst & med. & worst & med. & worst & med. & worst\\
\midrule
monodisperse       & 0.181 & 16.79 & 0.212 & 1.48 & 0.099 & 0.76 & 0.568 & 0.96\\
decoupling         & 0.558 & 286.07 & 0.716 & 28.22 & 0.167 & 5.92 & \textbf{0.287} & \textbf{0.42}\\
local monodisperse & \textbf{0.059} & 17.06 & \textbf{0.071} & 1.50 & 0.068 & 0.76 & 0.566 & 0.95\\
partial $S_{ij}$   & 0.558 & 286.11 & 0.717 & 28.22 & 0.167 & 5.92 & 0.290 & \textbf{0.42}\\
scaling            & 0.072 & 66.42 & 0.129 & 6.36 & 0.059 & 1.71 & 0.417 & 0.79\\
vdW1               & 0.085 & \textbf{11.87} & 0.088 & 2.00 & \textbf{0.044} & \textbf{0.61} & 0.575 & 1.11\\
\bottomrule
\end{tabular}
\end{table}

\subsubsection{The decoupling approximation}

The decoupling and partial-structure-factor schemes track each other closely
throughout, which is expected: both replace the distinct pair correlations by
a single common structure factor. Measured against the best scheme available
at the same state point, the decoupling approximation is typically worse by a
factor of

\begin{center}\small
\begin{tabular}{lcc}
\toprule
System & median ratio & worst\\
\midrule
hard sphere & 13.0 & 48.4\\
square well & 14.8 & 43.2\\
adhesive HS & 4.9 & 10.1\\
Yukawa & 1.0 & 1.6\\
\bottomrule
\end{tabular}
\end{center}

\noindent so its penalty falls monotonically as the interaction becomes softer
and longer ranged, until for the Yukawa fluid it disappears.

This is consistent with the known reason for its failure, and with an
independent result. Gazzillo \emph{et al.} (1999) reached the same conclusion
for polydisperse ionic fluids --- that the extended decoupling approximation
behaves poorly across essentially the whole $Q$ range while the scaling
approximation is accurate over the experimentally accessible range --- and
identified the cause: the decoupling approximation neglects the excluded-volume
conditions $g_{\alpha\beta}(r) = 0$ inside the hard cores, because it
replaces the distinct pair correlation functions by a small number of
effective ones. Where the structure factor is dominated by excluded volume the
approximation that discards excluded-volume detail loses most; where it is
dominated by a long-ranged soft tail, it loses least. Our implementation of
the scaling approximation retains both ingredients Gazzillo identifies as
necessary, the volume prefactor and evaluation of the reference structure
factor at a scaled argument, while our decoupling approximation uses a single
common structure factor, the $\lambda = 1$ limit in their notation.

\begin{quote}\small\itshape
Draft note. This section has now been written five times, on 4, 9, 60, 98 and
120 state points, and each version supported a different conclusion: that the
scheme ordering reverses between potentials (an artefact of the defect
described in \S\ref{sec:validation}); that the local monodisperse
approximation is uniformly best; that three schemes divide the plane; that the
systems fall into three regimes including one where the choice barely matters;
and the account above. The fourth of those was an artefact of having run only
the eight lowest-polydispersity adhesive points, which are indeed insensitive
to the choice --- the completed row is not. The $s = \varphi = 0.40$ corner
has now been checked and IS physical --- the diagonal partial structure
factors stay near $0.97$, the intensity is positive throughout, the packing
fraction reproduces exactly, and $\varphi = 0.40$ with a twofold spread in
diameter is nowhere near jamming, polydispersity raising the close-packing
threshold rather than lowering it. The factor of 286 therefore stands as a
genuine failure of the approximation at an ordinary state point, not a
numerical artefact.
\end{quote}

\subsection{Two cautions for the practitioner}

Both of the following were established while validating this work and are easy
to fall foul of.

\paragraph{Grid resolution is set by the narrowest feature.}
For potentials with a narrow attractive well the discretisation must resolve
the \emph{well}, not the particle: $\gtrsim 30\,\sigma/\delta$ points per
diameter. At the default resolution an adhesive-sphere calculation with
$\delta = 0.02\sigma$ is in error by ninety per cent, and the resulting
structure factor is smooth, plausible and wrong. In an earlier version of this
comparison this made all six schemes appear uniformly poor ($0.72$--$0.80$)
for that system, because each was being compared against a reference that was
itself badly in error.

\paragraph{Convergence flags and residuals are not sufficient.}
The closure equations admit multiple fixed points. At one strongly coupled
state two different accelerated solvers converged to residuals of
$\sim10^{-12}$ and returned $S(Q)=17.0$ and $9.7$ respectively, both with
$\min S(Q) < 0$. Since a structure factor is a variance it cannot be negative,
and solutions are therefore screened for physicality in addition to the
residual test. Separately, the SUNDIALS KINSOL solver was observed to report
success after diverging, because exponent clipping in the closure pins values
at $\sim10^{217}$ and defeats its internal stopping criterion; every solve
consequently recomputes its own residual before the result is accepted.
.
% Not standalone: no preamble, no \begin{document}.
%
% Numbers are the COMPLETE survey: 120 state points, 4 systems x 6
% polydispersities x 5 volume fractions, every point converged and screened.
% Source data: tools/survey.jsonl. Regenerate the tables and figure with
%     python tools/section53_survey.py --summarise survey.jsonl
%     python tools/section53_survey.py --plot survey.jsonl
%
% CAUTION FOR WHOEVER EDITS THIS. Five successive versions of this section
% have supported five different conclusions, each overturned by adding data.
% Do not tighten the language below without adding state points first.

\subsection{Comparison of the approximation schemes}
\label{sec:schemes}

We computed all six schemes alongside the exact multicomponent result over a
grid of relative polydispersity $s = 0.05$--$0.40$ and volume fraction
$\varphi = 0.05$--$0.40$ for four systems: hard spheres and adhesive hard
spheres ($\tau = 0.3$, $\delta = 0.02\sigma$) in the Percus--Yevick closure,
and a square-\emph{shoulder} fluid ($\varepsilon = +1\,k_BT$,
$\delta = 0.1\sigma$) and a
hard-core Yukawa fluid (screening length $0.5\sigma$, contact strength
$2\,k_BT$) in the hypernetted-chain closure. That is 120 state points, every
one converged and screened for physicality. The error measure is
$\max|I_{\mathrm{approx}}/I_{\mathrm{exact}} - 1|$ over
$0.3 \le Q\sigma \le 20$, with five structure-factor classes and sixty
form-factor classes. The narrow-well systems were computed at the resolution
their well width demands (\S\ref{sec:welltrap}): 300 points per diameter for
the square shoulder, 1500 for the adhesive spheres.

\subsubsection{Three systems share a pattern; the Yukawa fluid does not}

\begin{quote}\small\itshape
Note on the sign convention. The implementation writes the pair potential as
$u(r) = \varepsilon$ for $\sigma < r < \sigma+\delta$ and forms
$\exp(-u)$, so a POSITIVE $\varepsilon$ gives $\exp(-\varepsilon) < 1$ and is
REPULSIVE --- a square shoulder. Attraction requires $\varepsilon < 0$. The
survey was run at $\varepsilon = +1\,k_BT$ and is therefore a shoulder
throughout; the calculations are unaffected, but every description of this
row as a ``square well'' was wrong and has been corrected. Before submission,
consider adding a genuinely attractive row ($\varepsilon < 0$): an attractive
well and a repulsive shoulder are different physics, and the adhesive-sphere
row is the only attractive system currently in the survey.
\end{quote}

For the hard-sphere, square-shoulder and adhesive-sphere fluids the same three
schemes divide the $(s,\varphi)$ plane between them, in a consistent
arrangement (Table~\ref{tab:map}): the scaling approximation at low
polydispersity and low volume fraction, the local monodisperse approximation
through the middle of the plane and increasingly as $s$ grows, and the van der
Waals one-fluid scheme along the high-$\varphi$ edge. The boundaries move
between potentials --- the scaling region shrinks as attraction is added, and
vanishes altogether for the adhesive spheres --- but the arrangement does not.
In all three, the decoupling and partial-structure-factor schemes are never
the best choice at any state point.

The hard-core Yukawa fluid behaves differently in both respects. There the
decoupling approximation is the \emph{best} scheme at 21 of 30 points, and
every scheme carries a median error between 29 and 58\,\%. That is a statement
about which approximation is least bad rather than about any of them being
usable: a parameter extracted through any of these schemes for such a system
carries a systematic error of tens of per cent.

This is not an artefact of the particular Yukawa parameters. Varying the
screening length over a fivefold range and the contact strength by a factor
of two and a half leaves both halves of the statement intact: decoupling is
the best scheme at every state point tested, and the median error across all
six schemes stays between $0.55$ and $0.79$ throughout
(Table~\ref{tab:yukawa}).

\begin{table}[htbp]
\centering
\caption{The Yukawa behaviour is insensitive to the interaction parameters.
$\lambda$ is the screening length in units of $\sigma$ and $K$ the contact
strength in $k_BT$; the median is taken across all six approximation schemes.}
\label{tab:yukawa}
\small
\begin{tabular}{ccccll}
\toprule
$\lambda$ & $K$ & $s$ & $\varphi$ & best scheme & median\\
\midrule
0.5 & 2.0 & 0.20 & 0.20 & decoupling (0.42) & 0.60\\
0.5 & 2.0 & 0.30 & 0.30 & decoupling (0.35) & 0.79\\
0.2 & 2.0 & 0.20 & 0.20 & decoupling (0.18) & 0.55\\
0.2 & 2.0 & 0.30 & 0.30 & decoupling (0.24) & 0.76\\
1.0 & 2.0 & 0.20 & 0.20 & decoupling (0.39) & 0.70\\
0.5 & 5.0 & 0.20 & 0.20 & decoupling (0.44) & 0.73\\
\bottomrule
\end{tabular}
\end{table}

Two further combinations --- the longest-ranged and the strongest attraction,
both at $s = \varphi = 0.30$ --- failed to converge. That is expected on
approaching a spinodal rather than a defect, but it does mean the Yukawa
behaviour is characterised on its convergent side only.

\begin{table}[htbp]
\centering
\caption{Scheme with the smallest maximum relative error at each state point.
Three systems share an arrangement; the Yukawa fluid is the exception.
``dec'' is the decoupling and ``part'' the partial-$S_{ij}$ scheme.}
\label{tab:map}
\small
\begin{tabular}{c|ccccc|ccccc}
\toprule
& \multicolumn{5}{c|}{hard sphere, $\varphi$}
& \multicolumn{5}{c}{square shoulder, $\varphi$}\\
$s$ & 0.05 & 0.10 & 0.20 & 0.30 & 0.40 & 0.05 & 0.10 & 0.20 & 0.30 & 0.40\\
\midrule
0.05 & scal & scal & scal & scal & scal & scal & vdW1 & vdW1 & vdW1 & vdW1\\
0.10 & scal & scal & scal & vdW1 & vdW1 & scal & vdW1 & vdW1 & vdW1 & vdW1\\
0.15 & scal & scal & LMA  & LMA  & vdW1 & scal & LMA  & LMA  & LMA  & vdW1\\
0.20 & scal & LMA  & LMA  & LMA  & vdW1 & LMA  & LMA  & LMA  & LMA  & vdW1\\
0.30 & scal & LMA  & LMA  & LMA  & vdW1 & LMA  & LMA  & LMA  & LMA  & vdW1\\
0.40 & LMA  & LMA  & LMA  & LMA  & vdW1 & LMA  & LMA  & LMA  & LMA  & LMA\\
\midrule
& \multicolumn{5}{c|}{adhesive HS, $\varphi$}
& \multicolumn{5}{c}{Yukawa, $\varphi$}\\
$s$ & 0.05 & 0.10 & 0.20 & 0.30 & 0.40 & 0.05 & 0.10 & 0.20 & 0.30 & 0.40\\
\midrule
0.05 & vdW1 & vdW1 & vdW1 & vdW1 & vdW1 & scal & scal & dec  & dec  & dec\\
0.10 & vdW1 & vdW1 & vdW1 & vdW1 & vdW1 & scal & scal & dec  & dec  & dec\\
0.15 & vdW1 & vdW1 & vdW1 & vdW1 & vdW1 & scal & scal & dec  & dec  & dec\\
0.20 & vdW1 & vdW1 & LMA  & LMA  & vdW1 & dec  & scal & dec  & dec  & dec\\
0.30 & LMA  & LMA  & LMA  & LMA  & vdW1 & dec  & dec  & dec  & dec  & dec\\
0.40 & LMA  & LMA  & LMA  & LMA  & vdW1 & dec  & dec  & dec  & part & part\\
\bottomrule
\end{tabular}
\end{table}

\subsubsection{Typical and worst-case accuracy}

Over all 120 points the local monodisperse approximation and the van der Waals
one-fluid scheme are the best choice about equally often (38 and 37 points
respectively), the scaling approximation at 22 and the decoupling
approximation at 21 --- the latter almost entirely from the Yukawa system. The
partial-structure-factor scheme wins twice. Within-5\,\% rates over all
points are 35.8\,\% for vdW1, 31.7\,\% for both the local monodisperse and
scaling approximations, 16.7\,\% for the monodisperse approximation and
11.7\,\% for decoupling and partial $S_{ij}$.

Typical and worst-case accuracy do not agree, and the choice between them
depends on the application (Table~\ref{tab:schemestats}). For hard spheres the
local monodisperse approximation has the smallest median error but the van der
Waals one-fluid scheme the smallest worst case, by a factor of nearly six.
Reporting only one of these would mislead.

\begin{table}[htbp]
\centering
\caption{Median and worst-case error over the thirty state points of each
system. Best in each column in bold.}
\label{tab:schemestats}
\small
\begin{tabular}{lcccccccc}
\toprule
& \multicolumn{2}{c}{hard sphere} & \multicolumn{2}{c}{square well}
& \multicolumn{2}{c}{adhesive HS} & \multicolumn{2}{c}{Yukawa}\\
Scheme & med. & worst & med. & worst & med. & worst & med. & worst\\
\midrule
monodisperse       & 0.181 & 16.79 & 0.212 & 1.48 & 0.099 & 0.76 & 0.568 & 0.96\\
decoupling         & 0.558 & 286.07 & 0.716 & 28.22 & 0.167 & 5.92 & \textbf{0.287} & \textbf{0.42}\\
local monodisperse & \textbf{0.059} & 17.06 & \textbf{0.071} & 1.50 & 0.068 & 0.76 & 0.566 & 0.95\\
partial $S_{ij}$   & 0.558 & 286.11 & 0.717 & 28.22 & 0.167 & 5.92 & 0.290 & \textbf{0.42}\\
scaling            & 0.072 & 66.42 & 0.129 & 6.36 & 0.059 & 1.71 & 0.417 & 0.79\\
vdW1               & 0.085 & \textbf{11.87} & 0.088 & 2.00 & \textbf{0.044} & \textbf{0.61} & 0.575 & 1.11\\
\bottomrule
\end{tabular}
\end{table}

\subsubsection{The decoupling approximation}

The decoupling and partial-structure-factor schemes track each other closely
throughout, which is expected: both replace the distinct pair correlations by
a single common structure factor. Measured against the best scheme available
at the same state point, the decoupling approximation is typically worse by a
factor of

\begin{center}\small
\begin{tabular}{lcc}
\toprule
System & median ratio & worst\\
\midrule
hard sphere & 13.0 & 48.4\\
square well & 14.8 & 43.2\\
adhesive HS & 4.9 & 10.1\\
Yukawa & 1.0 & 1.6\\
\bottomrule
\end{tabular}
\end{center}

\noindent so its penalty falls monotonically as the interaction becomes softer
and longer ranged, until for the Yukawa fluid it disappears.

This is consistent with the known reason for its failure, and with an
independent result. Gazzillo \emph{et al.} (1999) reached the same conclusion
for polydisperse ionic fluids --- that the extended decoupling approximation
behaves poorly across essentially the whole $Q$ range while the scaling
approximation is accurate over the experimentally accessible range --- and
identified the cause: the decoupling approximation neglects the excluded-volume
conditions $g_{\alpha\beta}(r) = 0$ inside the hard cores, because it
replaces the distinct pair correlation functions by a small number of
effective ones. Where the structure factor is dominated by excluded volume the
approximation that discards excluded-volume detail loses most; where it is
dominated by a long-ranged soft tail, it loses least. Our implementation of
the scaling approximation retains both ingredients Gazzillo identifies as
necessary, the volume prefactor and evaluation of the reference structure
factor at a scaled argument, while our decoupling approximation uses a single
common structure factor, the $\lambda = 1$ limit in their notation.

\begin{quote}\small\itshape
Draft note. This section has now been written five times, on 4, 9, 60, 98 and
120 state points, and each version supported a different conclusion: that the
scheme ordering reverses between potentials (an artefact of the defect
described in \S\ref{sec:validation}); that the local monodisperse
approximation is uniformly best; that three schemes divide the plane; that the
systems fall into three regimes including one where the choice barely matters;
and the account above. The fourth of those was an artefact of having run only
the eight lowest-polydispersity adhesive points, which are indeed insensitive
to the choice --- the completed row is not. The $s = \varphi = 0.40$ corner
has now been checked and IS physical --- the diagonal partial structure
factors stay near $0.97$, the intensity is positive throughout, the packing
fraction reproduces exactly, and $\varphi = 0.40$ with a twofold spread in
diameter is nowhere near jamming, polydispersity raising the close-packing
threshold rather than lowering it. The factor of 286 therefore stands as a
genuine failure of the approximation at an ordinary state point, not a
numerical artefact.
\end{quote}

\subsection{Two cautions for the practitioner}

Both of the following were established while validating this work and are easy
to fall foul of.

\paragraph{Grid resolution is set by the narrowest feature.}
For potentials with a narrow attractive well the discretisation must resolve
the \emph{well}, not the particle: $\gtrsim 30\,\sigma/\delta$ points per
diameter. At the default resolution an adhesive-sphere calculation with
$\delta = 0.02\sigma$ is in error by ninety per cent, and the resulting
structure factor is smooth, plausible and wrong. In an earlier version of this
comparison this made all six schemes appear uniformly poor ($0.72$--$0.80$)
for that system, because each was being compared against a reference that was
itself badly in error.

\paragraph{Convergence flags and residuals are not sufficient.}
The closure equations admit multiple fixed points. At one strongly coupled
state two different accelerated solvers converged to residuals of
$\sim10^{-12}$ and returned $S(Q)=17.0$ and $9.7$ respectively, both with
$\min S(Q) < 0$. Since a structure factor is a variance it cannot be negative,
and solutions are therefore screened for physicality in addition to the
residual test. Separately, the SUNDIALS KINSOL solver was observed to report
success after diverging, because exponent clipping in the closure pins values
at $\sim10^{217}$ and defeats its internal stopping criterion; every solve
consequently recomputes its own residual before the result is accepted.
.
% Not standalone: no preamble, no \begin{document}.
%
% Numbers are the COMPLETE survey: 120 state points, 4 systems x 6
% polydispersities x 5 volume fractions, every point converged and screened.
% Source data: tools/survey.jsonl. Regenerate the tables and figure with
%     python tools/section53_survey.py --summarise survey.jsonl
%     python tools/section53_survey.py --plot survey.jsonl
%
% CAUTION FOR WHOEVER EDITS THIS. Five successive versions of this section
% have supported five different conclusions, each overturned by adding data.
% Do not tighten the language below without adding state points first.

\subsection{Comparison of the approximation schemes}
\label{sec:schemes}

We computed all six schemes alongside the exact multicomponent result over a
grid of relative polydispersity $s = 0.05$--$0.40$ and volume fraction
$\varphi = 0.05$--$0.40$ for four systems: hard spheres and adhesive hard
spheres ($\tau = 0.3$, $\delta = 0.02\sigma$) in the Percus--Yevick closure,
and a square-\emph{shoulder} fluid ($\varepsilon = +1\,k_BT$,
$\delta = 0.1\sigma$) and a
hard-core Yukawa fluid (screening length $0.5\sigma$, contact strength
$2\,k_BT$) in the hypernetted-chain closure. That is 120 state points, every
one converged and screened for physicality. The error measure is
$\max|I_{\mathrm{approx}}/I_{\mathrm{exact}} - 1|$ over
$0.3 \le Q\sigma \le 20$, with five structure-factor classes and sixty
form-factor classes. The narrow-well systems were computed at the resolution
their well width demands (\S\ref{sec:welltrap}): 300 points per diameter for
the square shoulder, 1500 for the adhesive spheres.

\subsubsection{Three systems share a pattern; the Yukawa fluid does not}

\begin{quote}\small\itshape
Note on the sign convention. The implementation writes the pair potential as
$u(r) = \varepsilon$ for $\sigma < r < \sigma+\delta$ and forms
$\exp(-u)$, so a POSITIVE $\varepsilon$ gives $\exp(-\varepsilon) < 1$ and is
REPULSIVE --- a square shoulder. Attraction requires $\varepsilon < 0$. The
survey was run at $\varepsilon = +1\,k_BT$ and is therefore a shoulder
throughout; the calculations are unaffected, but every description of this
row as a ``square well'' was wrong and has been corrected. Before submission,
consider adding a genuinely attractive row ($\varepsilon < 0$): an attractive
well and a repulsive shoulder are different physics, and the adhesive-sphere
row is the only attractive system currently in the survey.
\end{quote}

For the hard-sphere, square-shoulder and adhesive-sphere fluids the same three
schemes divide the $(s,\varphi)$ plane between them, in a consistent
arrangement (Table~\ref{tab:map}): the scaling approximation at low
polydispersity and low volume fraction, the local monodisperse approximation
through the middle of the plane and increasingly as $s$ grows, and the van der
Waals one-fluid scheme along the high-$\varphi$ edge. The boundaries move
between potentials --- the scaling region shrinks as attraction is added, and
vanishes altogether for the adhesive spheres --- but the arrangement does not.
In all three, the decoupling and partial-structure-factor schemes are never
the best choice at any state point.

The hard-core Yukawa fluid behaves differently in both respects. There the
decoupling approximation is the \emph{best} scheme at 21 of 30 points, and
every scheme carries a median error between 29 and 58\,\%. That is a statement
about which approximation is least bad rather than about any of them being
usable: a parameter extracted through any of these schemes for such a system
carries a systematic error of tens of per cent.

This is not an artefact of the particular Yukawa parameters. Varying the
screening length over a fivefold range and the contact strength by a factor
of two and a half leaves both halves of the statement intact: decoupling is
the best scheme at every state point tested, and the median error across all
six schemes stays between $0.55$ and $0.79$ throughout
(Table~\ref{tab:yukawa}).

\begin{table}[htbp]
\centering
\caption{The Yukawa behaviour is insensitive to the interaction parameters.
$\lambda$ is the screening length in units of $\sigma$ and $K$ the contact
strength in $k_BT$; the median is taken across all six approximation schemes.}
\label{tab:yukawa}
\small
\begin{tabular}{ccccll}
\toprule
$\lambda$ & $K$ & $s$ & $\varphi$ & best scheme & median\\
\midrule
0.5 & 2.0 & 0.20 & 0.20 & decoupling (0.42) & 0.60\\
0.5 & 2.0 & 0.30 & 0.30 & decoupling (0.35) & 0.79\\
0.2 & 2.0 & 0.20 & 0.20 & decoupling (0.18) & 0.55\\
0.2 & 2.0 & 0.30 & 0.30 & decoupling (0.24) & 0.76\\
1.0 & 2.0 & 0.20 & 0.20 & decoupling (0.39) & 0.70\\
0.5 & 5.0 & 0.20 & 0.20 & decoupling (0.44) & 0.73\\
\bottomrule
\end{tabular}
\end{table}

Two further combinations --- the longest-ranged and the strongest attraction,
both at $s = \varphi = 0.30$ --- failed to converge. That is expected on
approaching a spinodal rather than a defect, but it does mean the Yukawa
behaviour is characterised on its convergent side only.

\begin{table}[htbp]
\centering
\caption{Scheme with the smallest maximum relative error at each state point.
Three systems share an arrangement; the Yukawa fluid is the exception.
``dec'' is the decoupling and ``part'' the partial-$S_{ij}$ scheme.}
\label{tab:map}
\small
\begin{tabular}{c|ccccc|ccccc}
\toprule
& \multicolumn{5}{c|}{hard sphere, $\varphi$}
& \multicolumn{5}{c}{square shoulder, $\varphi$}\\
$s$ & 0.05 & 0.10 & 0.20 & 0.30 & 0.40 & 0.05 & 0.10 & 0.20 & 0.30 & 0.40\\
\midrule
0.05 & scal & scal & scal & scal & scal & scal & vdW1 & vdW1 & vdW1 & vdW1\\
0.10 & scal & scal & scal & vdW1 & vdW1 & scal & vdW1 & vdW1 & vdW1 & vdW1\\
0.15 & scal & scal & LMA  & LMA  & vdW1 & scal & LMA  & LMA  & LMA  & vdW1\\
0.20 & scal & LMA  & LMA  & LMA  & vdW1 & LMA  & LMA  & LMA  & LMA  & vdW1\\
0.30 & scal & LMA  & LMA  & LMA  & vdW1 & LMA  & LMA  & LMA  & LMA  & vdW1\\
0.40 & LMA  & LMA  & LMA  & LMA  & vdW1 & LMA  & LMA  & LMA  & LMA  & LMA\\
\midrule
& \multicolumn{5}{c|}{adhesive HS, $\varphi$}
& \multicolumn{5}{c}{Yukawa, $\varphi$}\\
$s$ & 0.05 & 0.10 & 0.20 & 0.30 & 0.40 & 0.05 & 0.10 & 0.20 & 0.30 & 0.40\\
\midrule
0.05 & vdW1 & vdW1 & vdW1 & vdW1 & vdW1 & scal & scal & dec  & dec  & dec\\
0.10 & vdW1 & vdW1 & vdW1 & vdW1 & vdW1 & scal & scal & dec  & dec  & dec\\
0.15 & vdW1 & vdW1 & vdW1 & vdW1 & vdW1 & scal & scal & dec  & dec  & dec\\
0.20 & vdW1 & vdW1 & LMA  & LMA  & vdW1 & dec  & scal & dec  & dec  & dec\\
0.30 & LMA  & LMA  & LMA  & LMA  & vdW1 & dec  & dec  & dec  & dec  & dec\\
0.40 & LMA  & LMA  & LMA  & LMA  & vdW1 & dec  & dec  & dec  & part & part\\
\bottomrule
\end{tabular}
\end{table}

\subsubsection{Typical and worst-case accuracy}

Over all 120 points the local monodisperse approximation and the van der Waals
one-fluid scheme are the best choice about equally often (38 and 37 points
respectively), the scaling approximation at 22 and the decoupling
approximation at 21 --- the latter almost entirely from the Yukawa system. The
partial-structure-factor scheme wins twice. Within-5\,\% rates over all
points are 35.8\,\% for vdW1, 31.7\,\% for both the local monodisperse and
scaling approximations, 16.7\,\% for the monodisperse approximation and
11.7\,\% for decoupling and partial $S_{ij}$.

Typical and worst-case accuracy do not agree, and the choice between them
depends on the application (Table~\ref{tab:schemestats}). For hard spheres the
local monodisperse approximation has the smallest median error but the van der
Waals one-fluid scheme the smallest worst case, by a factor of nearly six.
Reporting only one of these would mislead.

\begin{table}[htbp]
\centering
\caption{Median and worst-case error over the thirty state points of each
system. Best in each column in bold.}
\label{tab:schemestats}
\small
\begin{tabular}{lcccccccc}
\toprule
& \multicolumn{2}{c}{hard sphere} & \multicolumn{2}{c}{square well}
& \multicolumn{2}{c}{adhesive HS} & \multicolumn{2}{c}{Yukawa}\\
Scheme & med. & worst & med. & worst & med. & worst & med. & worst\\
\midrule
monodisperse       & 0.181 & 16.79 & 0.212 & 1.48 & 0.099 & 0.76 & 0.568 & 0.96\\
decoupling         & 0.558 & 286.07 & 0.716 & 28.22 & 0.167 & 5.92 & \textbf{0.287} & \textbf{0.42}\\
local monodisperse & \textbf{0.059} & 17.06 & \textbf{0.071} & 1.50 & 0.068 & 0.76 & 0.566 & 0.95\\
partial $S_{ij}$   & 0.558 & 286.11 & 0.717 & 28.22 & 0.167 & 5.92 & 0.290 & \textbf{0.42}\\
scaling            & 0.072 & 66.42 & 0.129 & 6.36 & 0.059 & 1.71 & 0.417 & 0.79\\
vdW1               & 0.085 & \textbf{11.87} & 0.088 & 2.00 & \textbf{0.044} & \textbf{0.61} & 0.575 & 1.11\\
\bottomrule
\end{tabular}
\end{table}

\subsubsection{The decoupling approximation}

The decoupling and partial-structure-factor schemes track each other closely
throughout, which is expected: both replace the distinct pair correlations by
a single common structure factor. Measured against the best scheme available
at the same state point, the decoupling approximation is typically worse by a
factor of

\begin{center}\small
\begin{tabular}{lcc}
\toprule
System & median ratio & worst\\
\midrule
hard sphere & 13.0 & 48.4\\
square well & 14.8 & 43.2\\
adhesive HS & 4.9 & 10.1\\
Yukawa & 1.0 & 1.6\\
\bottomrule
\end{tabular}
\end{center}

\noindent so its penalty falls monotonically as the interaction becomes softer
and longer ranged, until for the Yukawa fluid it disappears.

This is consistent with the known reason for its failure, and with an
independent result. Gazzillo \emph{et al.} (1999) reached the same conclusion
for polydisperse ionic fluids --- that the extended decoupling approximation
behaves poorly across essentially the whole $Q$ range while the scaling
approximation is accurate over the experimentally accessible range --- and
identified the cause: the decoupling approximation neglects the excluded-volume
conditions $g_{\alpha\beta}(r) = 0$ inside the hard cores, because it
replaces the distinct pair correlation functions by a small number of
effective ones. Where the structure factor is dominated by excluded volume the
approximation that discards excluded-volume detail loses most; where it is
dominated by a long-ranged soft tail, it loses least. Our implementation of
the scaling approximation retains both ingredients Gazzillo identifies as
necessary, the volume prefactor and evaluation of the reference structure
factor at a scaled argument, while our decoupling approximation uses a single
common structure factor, the $\lambda = 1$ limit in their notation.

\begin{quote}\small\itshape
Draft note. This section has now been written five times, on 4, 9, 60, 98 and
120 state points, and each version supported a different conclusion: that the
scheme ordering reverses between potentials (an artefact of the defect
described in \S\ref{sec:validation}); that the local monodisperse
approximation is uniformly best; that three schemes divide the plane; that the
systems fall into three regimes including one where the choice barely matters;
and the account above. The fourth of those was an artefact of having run only
the eight lowest-polydispersity adhesive points, which are indeed insensitive
to the choice --- the completed row is not. The $s = \varphi = 0.40$ corner
has now been checked and IS physical --- the diagonal partial structure
factors stay near $0.97$, the intensity is positive throughout, the packing
fraction reproduces exactly, and $\varphi = 0.40$ with a twofold spread in
diameter is nowhere near jamming, polydispersity raising the close-packing
threshold rather than lowering it. The factor of 286 therefore stands as a
genuine failure of the approximation at an ordinary state point, not a
numerical artefact.
\end{quote}

\subsection{Two cautions for the practitioner}

Both of the following were established while validating this work and are easy
to fall foul of.

\paragraph{Grid resolution is set by the narrowest feature.}
For potentials with a narrow attractive well the discretisation must resolve
the \emph{well}, not the particle: $\gtrsim 30\,\sigma/\delta$ points per
diameter. At the default resolution an adhesive-sphere calculation with
$\delta = 0.02\sigma$ is in error by ninety per cent, and the resulting
structure factor is smooth, plausible and wrong. In an earlier version of this
comparison this made all six schemes appear uniformly poor ($0.72$--$0.80$)
for that system, because each was being compared against a reference that was
itself badly in error.

\paragraph{Convergence flags and residuals are not sufficient.}
The closure equations admit multiple fixed points. At one strongly coupled
state two different accelerated solvers converged to residuals of
$\sim10^{-12}$ and returned $S(Q)=17.0$ and $9.7$ respectively, both with
$\min S(Q) < 0$. Since a structure factor is a variance it cannot be negative,
and solutions are therefore screened for physicality in addition to the
residual test. Separately, the SUNDIALS KINSOL solver was observed to report
success after diverging, because exponent clipping in the closure pins values
at $\sim10^{217}$ and defeats its internal stopping criterion; every solve
consequently recomputes its own residual before the result is accepted.
